\section*{2. Same Data, Different Similarity}

令
\[
 A=(5,1,5,0),\quad B=(10,2,10,0),\quad
 C=(5,5,1,5),\quad D=(1,0,1,0).
\]

\begin{enumerate}
  \item $A$ 與其他學生的 Euclidean distance 為
  \begin{align*}
    d(A,B)&=\sqrt{(5-10)^2+(1-2)^2+(5-10)^2+(0-0)^2}
           =\sqrt{51}\approx 7.141,\\
    d(A,C)&=\sqrt{(5-5)^2+(1-5)^2+(5-1)^2+(0-5)^2}
           =\sqrt{57}\approx 7.550,\\
    d(A,D)&=\sqrt{(5-1)^2+(1-0)^2+(5-1)^2+(0-0)^2}
           =\sqrt{33}\approx 5.745.
  \end{align*}
  因此，由近至遠的排序為
  \[
    \boxed{D\;\rightarrow\;B\;\rightarrow\;C}.
  \]

  \item 由向量長度與內積可得
  \begin{align*}
    \cos(A,B)&=\frac{102}{\sqrt{51}\sqrt{204}}=1.0000,\\
    \cos(A,C)&=\frac{35}{\sqrt{51}\sqrt{76}}
              =\frac{35}{\sqrt{3876}}\approx 0.5621,\\
    \cos(A,D)&=\frac{10}{\sqrt{51}\sqrt{2}}
              =\frac{10}{\sqrt{102}}\approx 0.9901.
  \end{align*}
  因此，由最相似至最不相似的排序為
  \[
    \boxed{B\;\rightarrow\;D\;\rightarrow\;C}.
  \]

  \item 兩種度量找出的最相似學生不同。Euclidean distance 選擇 $D$，因為 $D$
  的絕對活動次數較小，在數值大小上離 $A$ 最近。Cosine similarity 則選擇 $B$；
  由於 $B=2A$，兩人的活動比例與向量方向完全相同，即使 $B$ 的所有非零次數都是
  $A$ 的兩倍。前者重視絕對位移，後者重視相對活動模式。

  \item 年收入可能以數萬為單位，遠大於四個活動次數。計算原始 Euclidean distance
  時，收入差的平方可能壓過其他屬性，使結果幾乎只反映收入。我會使用訓練資料擬合
  standardization，將所有量化屬性轉為 $z$-score。若收入分布高度偏斜或含有極端值，
  也可先做 log transformation，或改用 median 與 interquartile range 的 robust scaling。
\end{enumerate}

\clearpage
